% from ega4-errata-addenda-list3-fr.tex
\oldpage[IV]{352}
\paragraph{(Err\textsubscript{IV}, 27).}
Dans (IV, 7.8.4), ajouter, avant la ligne 6 du bas de la p.~217 :

\smallskip
\noindent
\textup{(vi)} Soit $A$ un anneau local noethérien intègre de dimension
1; il est alors universellement caténaire
{\hyperref[canon-index.ega4.c04.cor.7-2-9]{(7.2.9)}} et, pour qu'il soit excellent,
il faut et il suffit donc, en vertu de
{\hyperref[canon-index.ega4.c04.statement.7-3-19]{(7.3.19)}}, que ses fibres
formelles soient géométriquement réduites, ou encore que son complété
$\widehat A$ soit réduit et que les corps composants de son anneau total
de fractions soient des extensions séparables du corps des fractions
$K$ de $A$. En vertu de {\hyperref[canon-index.ega4.c04.thm.7-6-4]{(7.6.4)}} et
{\hyperref[canon-index.ega4.c04.thm.7-7-2]{(7.7.2)}}, il revient aussi au même de
dire que $A$ est universellement japonais. Lorsque $A$ est un anneau de
valuation discrète, il en est de même de $\widehat A$; dire que $A$ est
excellent signifie alors que le corps des fractions de $\widehat A$
est une extension séparable de $K$, condition toujours satisfaite si
$K$ est de caractéristique 0.

\paragraph{(Err\textsubscript{IV}, 28).}
Après (IV, 8.10.3), après la ligne 14 de la p.~37, insérer :

\begin{remark}[8.10.3.1]
\label{IV.8.10.3.1}
On ne peut dans {\hyperref[canon-index.ega4.c04.prop.8-10-3]{(8.10.3)}} supprimer
l'hypothèse que les $\overline{f_\alpha(X_\alpha)}$ soient
constructibles. Prenons par exemple $S=\operatorname{Spec}(A)$, où $A$
est un anneau tel que $S$ soit compact, totalement discontinu et sans
point isolé, tout idéal premier $\mathfrak p$ de $A$ étant maximal et
$A_\mathfrak p$ étant un corps (Bourbaki, \emph{Alg. comm.}, chap.~II,
\S4, exerc.~16 et 17). Soit $s\in S$, et prenons pour
$Y_\alpha=S_\alpha$ les voisinages ouverts affines de $s$ dans $S$; on
prend $X_\alpha=\operatorname{Spec}(k(s))$ pour tout $\alpha$,
$f_\alpha$ et $f$ étant les morphismes canoniques; il est clair que
$f$, qui est l'identité, est dominant, mais aucun des $f_\alpha$ ne
l'est.
\end{remark}

\paragraph{(Err\textsubscript{IV}, 29).}
Dans (IV, 8.12.8), ligne 7 de la p.~47, remplacer « tel qu'il existe un
$\sh{O}_Y$-Module ample » par : « et quasi-séparé ». Ligne 12 de la
p.~47, remplacer « L'hypothèse » par : « Il résulte de
{\hyperref[canon-index.ega4.c04.prop.8-12-3]{(8.12.3)}} que l'existence de la
factorisation de l'énoncé est une propriété locale sur $Y$; on peut
donc supposer $Y$ affine, ce qui ».

\paragraph{(Err\textsubscript{IV}, 30).}
Après (IV, 8.12.11), ajouter à la fin de la p.~48 :

\begin{corollary}[8.12.12]
\EGAReaderTarget{ega4.c04.cor.8-12-12}{IV.8.12.12}
Soient $X$, $Y$ deux préschémas intègres,
$f:X\to Y$ un morphisme séparé, birationnel et localement de type fini.
Soit $y\in Y$ un point normal de $Y$ et supposons qu'un point
$x\in f^{-1}(y)$ soit isolé dans $f^{-1}(y)$. Alors il existe un
voisinage ouvert $U$ de $y$ dans $Y$ tel que la restriction
$f^{-1}(U)\to U$ de $f$ soit un isomorphisme.
\end{corollary}

\begin{proof}
Il suffit de prouver qu'il existe un voisinage ouvert $U$ de $y$ dans
$Y$ et une $U$-section $g:U\to f^{-1}(U)$ de $f^{-1}(U)$ telle que
$g(y)=x$. En effet, comme $U$ contient le point générique $\eta$ de
$Y$ et $f^{-1}(U)$ le point générique $\xi$ de $X$, on a
nécessairement $g(\eta)=\xi$ puisque $f$ est birationnel
{\hyperref[canon-index.ega4.c04.statement.6-15-4]{(6.15.4)}}; comme $g$ est une
immersion, $g(U)$ est localement fermé dans $f^{-1}(U)$, donc ouvert
dans son adhérence dans $f^{-1}(U)$, et comme cette dernière est égale
à $f^{-1}(U)$ d'après ce qui précède, $g(U)$ est ouvert dans
$f^{-1}(U)$; comme le sous-préschéma de $f^{-1}(U)$ associé à $g$ est
isomorphe à $U$, donc réduit, et que $f^{-1}(U)$ est réduit, il résulte
de (I, 5.2.1) que $g$ est une immersion ouverte. D'autre part, puisque
$f$ est séparé, $g$ est une immersion fermée (I, 5.4.6), donc $g(U)$
est un ensemble à la fois ouvert et fermé dans $f^{-1}(U)$. Mais comme
$f^{-1}(U)$ est intègre, cela entraînera $g(U)=f^{-1}(U)$, d'où la
conclusion.

Considérons alors deux cas :

\textup{I)} Supposons d'abord $Y$ normal, et soit $V$ l'ensemble des
points de $X$ en lesquels $f$ est quasi-fini; c'est un ensemble ouvert
dans $X$ {\hyperref[canon-index.ega4.c04.cor.13-1-4]{(13.1.4)}}.\footnote{Le
lecteur vérifiera que ni {\hyperlink{ega4.c04.cor.8-12-12}{(8.12.12)}}
ni aucune de ses conséquences n'est utilisée dans la preuve de
{\hyperref[canon-index.ega4.c04.cor.13-1-4]{(13.1.4)}}.} Il résulte alors de
{\hyperref[canon-index.ega4.c04.cor.8-12-10]{(8.12.10)}}

